% ---------------------------------------------------
% Trabajo Final de Grado
% Author: Fabián González Lence <alu0101549491@ull.edu.es>
% Chapter: Goals 
% ----------------------------------------------------


\chapter{Objetivos} \label{chap:Goals}

El objetivo principal de este \ac{TFG} es evaluar de forma práctica las capacidades actuales de la \ac{IA} generativa para desarrollar aplicaciones web no triviales, analizando las prestaciones y rendimiento de varios \ac{LLMs} en las diferentes etapas del \ac{SDLC}.\\

\noindent Para alcanzar dicho objetivo general, se plantean los siguientes objetivos específicos:

\begin{enumerate}
  \item \textbf{Diseñar una metodología experimental basada en roles de IA.} Definir un flujo de trabajo en el que distintos modelos de \ac{IA} asuman roles especializados ---arquitecto, desarrollador, revisor de código y generador de pruebas--- y cooperen de forma coordinada a lo largo del
  Ciclo de Vida del Desarrollo de Software, (\textit{Software Development Life Cycle}, SDLC).

  \item \textbf{Desarrollar cinco aplicaciones de complejidad creciente.} Construir, mediante la colaboración entre modelos de IA las siguientes aplicaciones:
  \begin{itemize}
    \item \textbf{\ac{HANGMAN}}: Versión web del conocido juego del ahorcado implementado en Typescript como una \ac{SPA}, implementando el patrón MVC  \cite{Reenskaug:1979:MVC} y usando la API Canvas de HTML.
    \item \textbf{\ac{PLAYER}}: Reproductor de música web interactivo simple desarrolado con React\footnote{React: \url{https://react.dev/}}, TypeScript\footnote{TypeScript: \url{https://www.typescriptlang.org/}} y Vite\footnote{Vite: \url{https://vite.dev/}}. 
    \item \textbf{\ac{BALATRO}}: Juego de cartas inspirado en Balatro\footnote{Balatro \url{https://es.wikipedia.org/wiki/Balatro}} con sistema de progresión por niveles, interacciones con una tienda y utiliación de manos de póker para el cálculo de puntos.
    \item \textbf{\ac{CARTO}}: Sistema de gestión de proyectos cartográficos implementando Clean Architecture \cite{Martin:2018:CleanArchitecture}, Vue.js\footnote{Vue: \url{https://vuejs.org/}} y comunicación en tiempo real.
    \item \textbf{\ac{TENNIS}}: Aplicación de gestión integral de torneos de tenis con múltiples roles de usuario construida en Angular\footnote{Angular: \url{https://angular.dev/}}.
  \end{itemize}

  \item \textbf{Documentar el proceso de interacción con las IAs.} Documentar de forma sistemática los \textit{prompts} empleados, las respuestas obtenidas y las decisiones tomadas en cada fase del desarrollo, de modo que el proceso sea reproducible y analizable.

  \item \textbf{Aplicar un proceso sistemático de revisión y testing asistido por IA.} Someter el código generado a revisiones de calidad y a la generación automatizada de conjuntos de pruebas, evaluando la capacidad de los modelos para detectar defectos en código producido por otros modelos.

  \item \textbf{Evaluar la calidad del software resultante.} Analizar los proyectos desarrollados aplicando métricas de calidad de código ---funcionalidad, mantenibilidad, coherencia, cobertura de pruebas y adherencia a buenas prácticas---.

  \item \textbf{Extraer conclusiones sobre la viabilidad del enfoque colaborativo multimodelo.} Sintetizar los hallazgos obtenidos para determinar en qué medida la cooperación entre múltiples IAs especializadas puede sustituir o complementar el trabajo humano en el proceso de desarrollo de software, y proponer líneas futuras de investigación.
\end{enumerate}




\begin{comment}
Este capítulo se escribe normalmente al final de la realización del trabajo. 
En ese momento se tendrá una perspectiva más clara de lo que se ha pretendido hacer y de lo que realmente se ha hecho.

Los objetivos específicos a conseguir en el desarrollo de este software son los siguientes:

1.Creación del anteproyecto. Elaboración de un documento donde se expone una
introducción, antecedentes y estado actual, así como un listado de tareas a realizar
y el plan de trabajo establecido para la realización de este Trabajo Fin de Grado.

2.Creación y configuración de la estructura principal del proyecto. Creación de los di-
rectorios en los que se alojará todo el código, distribuido en los directorios frontend,
back-end y algorithm.

3.Creación de la API-REST. Desarrollo de un back-end capaz de comunicar la interfaz
de usuario con una base de datos como lo es MongoDB, haciendo uso del lenguaje
de programación Java y su framework Spring.

4.Creación de la interfaz de usuario. Desarrollo de un front-end desde el que el usuario
pueda interactuar con la aplicación. Para el desarrollo de esta interfaz se hizo uso
del framework Vue.js y Vuetify.

5.Implementación e integración en la aplicación del BAP. Implementación de un
módulo, aparte de los dos anteriores, capaz de resolver el roblema de gestión de
atraques (BAP).

6.Creación de una pantalla de resultados. Desarrollo de una pantalla en la que se
muestre de manera gráfica y detallada la salida obtenida al aplicar el BAP.

7.Redacción de la memoria. Se describe el Trabajo Fin de Grado desarrollado.


This document summarizes the research and development work carried out by the student in the achievement of his Final Degree Project (\textit{Trabajo de Fin de Grado}, TFG), which will conclude his studies for the degree \textit{Grado en Ingeniería Informática} at the \textit{Escuela Superior de Ingeniería y Tecnología} at the University of La Laguna (ULL).

This project has the following main goals:

\begin{enumerate}
    \item A first objective has been ... 
    \item Another key objective is ... 
    \item The implementation and ...
    \item Additionally, the student will apply....
    \item Finally, after thorough experimentation and evaluation, the student is expected to deliver...
\end{enumerate}
\end{comment}